\subsection{Graphing General Exponential Functions} 
We can also compose an exponential function with \makeblank[1in] functions. The most general form is
\[
f(x)=a(b)^{cx-h}+k
\]
To graph this, we will be explicit about a strategy we have been using for a while: looking at the input and output of the ``base function.''
\begin{example}
Graph $f(x)=-2(3)^{\frac{1}{2}x-3}+5$
\begin{cpic}{}{16}
	\draw (-5,-2) -- (5,-2);
	\foreach \x in {-3,0,3}
		\draw (\x,-1) -- (\x,-7);
	\foreach \x/\lbl in {-4/$x$,-1.5/input,1.5/output,4/$y$}
		\regtext{\x}{-2}{\lbl};
\end{cpic}
\end{example}
Here are the key features we will use to graph.
\begin{itemize}
\item The horizontal asymptote at \makeblank[1in].
\item If \makeblanksmall is negative, the graph is \makeblank[1in] the asymptote.
\item To find two key points, we will use the points where the \makeblank[1in] is \makeblanksmall, \makeblanksmall, and \makeblanksmall; these will usually give us at least two nice points and one fraction.
\end{itemize}